\section{Нефункциональные требования}

\subsection{Производительность}

Время загрузки и выполнения проверки одного PDF не должно превышать приемлемого порога (15\,с для типового документа). При пакетной обработке деградация времени отклика недопустима. Интерфейс должен оставаться отзывчивым при типичных сценариях использования.

\subsection{Безопасность}

Используется HTTPS для всего трафика. Аутентификация~--- предпочтительно itmo.id; применяется разграничение доступа по ролям. Все вводимые пользователем данные валидируются и санитизируются. Для загружаемых PDF обязательно выполняется проверка безопасности: не допускаются файлы с вредоносным ПО, исполняемым кодом или PDF-эксплойтами (встроенный JavaScript, опасные действия и т.\,п.).

\subsection{Удобство использования (Usability)}

Интерфейс на русском языке, с понятными сообщениями об ошибках и подсказками. Интерфейс адаптируется к разным размерам экрана (десктоп, планшет).

\subsection{Надёжность}

Обеспечивается логирование, корректная обработка ошибок и информативные сообщения пользователю. SLA 99.5\% в период учебной практики (допускаются техработы ночью).

\subsection{Поддерживаемость}

Код бэкенда и фронтенда следует принятым практикам (FastAPI, React). API документируется (OpenAPI). Архитектура модульная. Непрерывная интеграция и развёртывание описаны в разделе «Архитектура системы».
